\documentclass[11pt]{amsart}
\usepackage[margin=1.1in]{geometry}
\usepackage{amssymb}
\usepackage{booktabs}
\usepackage{array}
\usepackage[hidelinks]{hyperref}
\usepackage{xurl}

% Generated by verify.py --write-tables from data/; do not edit.
\newcommand{\TenTwo}{212}
\newcommand{\TenThree}{44}
\newcommand{\TenFive}{146}
\newcommand{\TenOrder}{373}
\newcommand{\TenEdges}{673}
\newcommand{\TenLeaves}{22}
\newcommand{\TenBlocks}{45}
\newcommand{\TenTypes}{20}
\newcommand{\TenGamma}{94}
\newcommand{\BlockRows}{%
1 & 3 & 4 & $2^{3}\cdot 3\cdot 5^{2}\cdot 11$ & $2^{2}\cdot 5\cdot 7\cdot 47$ \\
2 & 2 & 6 & $2^{2}\cdot 3\cdot 5^{4}\cdot 71$ & $2^{5}\cdot 5^{3}\cdot 7\cdot 19$ \\
3 & 1 & 6 & $2^{2}\cdot 5^{2}\cdot 73^{2}$ & $2^{3}\cdot 5^{3}\cdot 13\cdot 41$ \\
4 & 3 & 7 & $-2^{2}\cdot 3\cdot 5^{2}\cdot 19\cdot 29^{2}$ & $-2^{10}\cdot 5^{2}\cdot 11\cdot 17$ \\
5 & 4 & 7 & $-2^{2}\cdot 3\cdot 5^{3}\cdot 47\cdot 67$ & $-2\cdot 5\cdot 11^{2}\cdot 59\cdot 67$ \\
6 & 3 & 8 & $2^{4}\cdot 3\cdot 5^{3}\cdot 7^{2}\cdot 11\cdot 13$ & $2^{4}\cdot 5^{5}\cdot 29^{2}$ \\
7 & 5 & 8 & $2^{7}\cdot 3\cdot 5^{3}\cdot 17\cdot 53$ & $2^{11}\cdot 5^{3}\cdot 13^{2}$ \\
8 & 1 & 8 & $2^{6}\cdot 3\cdot 5^{2}\cdot 11\cdot 19\cdot 43$ & $2^{7}\cdot 5^{2}\cdot 7\cdot 41\cdot 47$ \\
9 & 2 & 8 & $2^{3}\cdot 3\cdot 5^{2}\cdot 17\cdot 43\cdot 97$ & $2^{7}\cdot 5^{2}\cdot 11\cdot 17\cdot 71$ \\
10 & 2 & 8 & $2^{7}\cdot 3\cdot 5^{3}\cdot 29\cdot 31$ & $2\cdot 5\cdot 13\cdot 17^{2}\cdot 31\cdot 37$ \\
11 & 1 & 8 & $2^{7}\cdot 3\cdot 5^{3}\cdot 29\cdot 31$ & $2^{2}\cdot 5^{3}\cdot 7\cdot 11\cdot 19\cdot 59$ \\
12 & 2 & 8 & $2^{2}\cdot 3\cdot 5^{2}\cdot 11\cdot 13\cdot 17\cdot 59$ & $2^{5}\cdot 5^{2}\cdot 11\cdot 67\cdot 73$ \\
13 & 2 & 8 & $2^{2}\cdot 3\cdot 5^{3}\cdot 23\cdot 29\cdot 43$ & $2^{5}\cdot 5^{3}\cdot 7\cdot 29\cdot 53$ \\
14 & 3 & 8 & $2^{2}\cdot 3\cdot 5^{2}\cdot 11\cdot 13\cdot 17\cdot 59$ & $2^{5}\cdot 5^{3}\cdot 7\cdot 29\cdot 53$ \\
15 & 1 & 8 & $2^{7}\cdot 3\cdot 5^{2}\cdot 67^{2}$ & $2\cdot 5\cdot 13\cdot 17^{2}\cdot 31\cdot 37$ \\
16 & 1 & 9 & $-2^{5}\cdot 3\cdot 5^{2}\cdot 7^{2}\cdot 37\cdot 89$ & $-2^{5}\cdot 5^{2}\cdot 7\cdot 23\cdot 31\cdot 97$ \\
17 & 1 & 9 & $-2^{7}\cdot 3\cdot 5^{3}\cdot 7\cdot 13\cdot 89$ & $-2^{5}\cdot 5^{2}\cdot 7\cdot 23\cdot 31\cdot 97$ \\
18 & 2 & 9 & $-2^{3}\cdot 3\cdot 5^{3}\cdot 31\cdot 37\cdot 113$ & $-2^{2}\cdot 5^{3}\cdot 7\cdot 11\cdot 89\cdot 113$ \\
19 & 1 & 9 & $-2^{3}\cdot 3\cdot 5^{3}\cdot 7\cdot 11\cdot 41^{2}$ & $-2^{2}\cdot 5^{3}\cdot 7^{4}\cdot 17\cdot 19$ \\
20 & 5 & 10 & $2^{8}\cdot 3\cdot 5^{4}\cdot 7^{2}\cdot 11\cdot 13$ & $2^{3}\cdot 5^{10}\cdot 43$ \\}
\newcommand{\BlockEdgeRows}{%
1 & $\{0,1\}$ & 01 03 12 23 \\
2 & $\{0,3,5\}$ & 01 04 12 15 23 34 35 45 \\
3 & $\{1,5\}$ & 01 02 04 05 12 15 23 25 34 \\
4 & $\{3,4\}$ & 01 05 12 16 23 24 26 34 45 56 \\
5 & $\{0,1,2,4,6\}$ & 01 02 04 05 14 15 23 24 25 34 35 36 \\
6 & $\{0\}$ & 01 04 07 12 14 16 23 27 45 47 67 \\
7 & $\{0,2\}$ & 01 04 07 12 23 25 37 45 56 67 \\
8 & $\{1,2,3,4,7\}$ & 01 03 04 12 17 23 24 25 26 37 47 57 67 \\
9 & $\{5,7\}$ & 01 02 07 12 14 16 17 23 34 37 45 57 \\
10 & $\{1,5,6\}$ & 01 04 07 12 14 25 27 36 37 45 47 56 57 \\
11 & $\{1,5,7\}$ & 01 04 07 12 14 25 27 36 37 45 47 56 57 \\
12 & $\{0,1,3,6\}$ & 01 02 07 12 17 23 26 34 35 36 45 47 57 67 \\
13 & $\{3,5\}$ & 01 05 07 12 14 17 23 24 27 36 37 45 46 57 67 \\
14 & $\{3\}$ & 01 02 04 05 14 17 23 25 27 36 37 47 56 57 67 \\
15 & $\{3,4\}$ & 02 03 05 07 12 13 17 23 24 27 34 45 46 57 67 \\
16 & $\{0,2,3,5,7\}$ & 01 02 03 04 05 12 13 15 18 24 25 27 28 34 37 46 47 56 68 \\
17 & $\{0,2,3,4,7,8\}$ & 01 02 05 13 18 24 26 36 37 45 56 67 68 \\
18 & $\{3,5,6,7\}$ & 01 02 03 13 14 15 24 25 36 37 48 68 78 \\
19 & $\{0,3,6,7,8\}$ & 01 02 06 07 13 14 24 35 36 46 57 58 67 78 \\
20 & $\{1,8\}$ & 01 03 06 08 12 14 24 26 28 35 46 49 67 \\}
\newcommand{\SixteenProd}{14929920}
\newcommand{\SixteenP}{58320}
\newcommand{\SixteenQ}{61440}
\newcommand{\SixteenEdges}{01,\ 02,\ 03,\ 13,\ 14,\ 19,\ 27,\ 37,\ 45,\ 46,\ 56,\ 58,\ 59,\ 67,\ 68,\ 69}
\newcommand{\SixteenT}{\{0,1,7,9\}}
\newcommand{\FourEdges}{$\{0,1\}$, $\{0,8\}$, $\{0,9\}$, $\{1,3\}$, $\{1,6\}$, $\{1,7\}$, $\{1,8\}$, $\{2,3\}$, $\{2,5\}$, $\{2,10\}$, $\{2,11\}$, $\{2,12\}$, $\{2,13\}$, $\{2,14\}$, $\{4,6\}$, $\{4,7\}$, $\{4,9\}$, $\{5,8\}$, $\{5,9\}$}
\newcommand{\SixteenPoly}{x^{16} + 16x^{15} + 115x^{14} + 498x^{13} + 1457x^{12} + 3042x^{11} + 4661x^{10} \\ + 5310x^{9} + 4508x^{8} + 2795x^{7} + 1215x^{6} + 346x^{5} + 57x^{4} + 4x^{3}}
\newcommand{\SixteenDeriv}{-35032640}
\newcommand{\FourDeriv}{13282240}
\newcommand{\ResidueRows}{%
$G_{4}$ & 15 & 0 & 930 & 768 & 281 & 773 \\
$G_{6}$ & 31 & 122 & 0 & 230 & 885 & 771 \\
$G_{8}$ & 83 & 230 & 101 & 0 & 298 & 620 \\
$G_{10}$ & 373 & 260 & 186 & 141 & 0 & 912 \\}


\newtheorem{theorem}{Theorem}
\newtheorem{corollary}[theorem]{Corollary}
\newtheorem{lemma}{Lemma}
\theoremstyle{remark}
\newtheorem{remark}{Remark}

\newcommand{\nmin}{n_{\min}}

\title[A domination root at $-10$]{A graph with domination root $-10$,\\ and small graphs with domination root $-4$}
\author{Aryan Poonacha}
\address{\url{https://github.com/Aryan-Poonacha/math-findings}}
\date{27 November 2026}
\subjclass[2020]{05C31, 05C69}
\keywords{domination polynomial, domination root, integer root}


\begin{document}
\begin{abstract}
The domination polynomial of a graph $G$ is $D(G,x)=\sum_S x^{|S|}$, summed over
the dominating sets $S$ of $G$. Alikhani and Griswold (2026) found a graph with
$D(G,-4)=0$, refuting the conjecture that $0$ and $-2$ are the only integer
domination roots, and asked about other roots and about the least order for
$-4$. We give a connected graph on $373$ vertices with $D(G,-10)=0$; with known
results, the integers in $[-11,-1]$ that are domination roots are exactly
$-2,-4,-6,-8,-10$. An exhaustive computer search shows that no graph on at most
$12$ vertices has domination root $-4$, so the least order is $13$, $14$ or
$15$. We also give an explicit $16$-vertex certificate with a simple root $-4$,
and extend the known joins so that $-4,-6,-8,-10$ occur together with any
prescribed multiplicities.
\end{abstract}
\maketitle

\section{Introduction}

All graphs are finite and simple. A set $S\subseteq V(G)$ \emph{dominates} $G$
if every vertex of $G$ lies in $S$ or has a neighbour in $S$. The
\emph{domination polynomial} is
\[
  D(G,x)=\sum_{S\subseteq V(G),\ S \text{ dominates } G} x^{|S|},
\]
and a \emph{domination root} of $G$ is a zero of $D(G,x)$. The polynomial is
monic of degree $|V(G)|$, and $D(G_1\cup G_2,x)=D(G_1,x)D(G_2,x)$ for a disjoint
union. For a nonempty graph, $0$ is a root, and $-2$ is a root for $K_2$, since
$D(K_2,x)=x^2+2x$. That $0$ and $-2$ are the only integer domination roots was
conjectured in Alikhani's 2009 thesis and in \cite{AAOP10}; see
\cite[Conjecture~1.1]{AG26}. Brouwer, Csorba and Schrijver \cite{BCS09} proved that every
graph has an odd number of dominating sets, that is, $D(G,1)$ is odd; since
$D(G,r)\equiv D(G,1)\pmod 2$ for every odd integer~$r$, odd integers are never
domination roots. Weisstein \cite{Wei} reports that no graph with at most $10$
vertices has an integer domination root other than $0$ and $-2$.

Alikhani and Griswold \cite{AG26} refuted the conjecture with a connected graph
$G_{33}$ on $33$ vertices with $D(G_{33},-4)=0$, and asked
\cite[Section~4]{AG26}:
\begin{quote}
1. Are there other negative integer domination roots, such as $-6$ or $-8$,
attainable via richer S-unit branch cancellations?\\
2. Is $n = 33$ the minimal order of a graph admitting $-4$ as a domination root?\\
3. Does there exist an infinite family of graphs $G$ satisfying $D(G,-4) = 0$?
\end{quote}
Several answers have appeared since, all as preprints, blog posts or
repositories:
\begin{itemize}
\item 18 September 2026: a connected planar graph with $30$ vertices and simple
  root $-4$ (Reshetnikov's repository \cite{Res26}), answering question~2
  negatively.
\item By 10 October 2026 (the date of the archived copy cited): the MathArena post
  \cite{MA26} gives a connected graph with $15$ vertices and $19$ edges with
  root $-4$ (found in a run of the GPT-6.1 Sol model) and connected graphs
  with roots $-6$ ($31$ vertices) and $-8$ ($83$ vertices). It reports an
  exhaustive search showing that no connected graph with at most $11$ vertices
  has root $-4$, so that $12\le \nmin(-4)\le 15$, and remarks that joining the
  centres of such graphs multiplies their domination polynomials, which gives
  each root at arbitrarily large orders and one graph with all three roots.
  MathArena also reports an accepted $16$-vertex witness from its first
  constructive run. Its published witness data specify the $15$-vertex graph.
  Below we give an explicit $16$-vertex certificate and verify that its root is
  simple.
\item 1 November 2026: the Omega Institute published a Lean-checked
  $21$-vertex witness for $-4$ \cite{Om26} and separately formalized a
  connected family of orders $33r+2$, with $r\ge1$ \cite{Om26b}.
\item 3 November 2026: Ma and Zhang \cite{MZ26} subsequently gave a written
  support-join proof of exact multiplicity $r$ and connected examples of every
  order $n\ge 35$, with maximum degree seven and cycle rank four.
\end{itemize}
The joining step rests on the irrelevant-edge theorem of Kotek, Preen, Simon,
Tittmann and Trinks \cite[Theorems~3.10 and~3.11]{KPSTT12} (Lemma~\ref{lem:join}
below). Here $\nmin(r)$ denotes the least order of a graph with domination
root~$r$; since $D$ is multiplicative over components, it is also the least
order of a connected one.

The main contributions are the $-10$ witness (Theorem~\ref{thm:ten}), which
with the results above determines the integer domination roots in $[-11,-1]$
(Corollary~\ref{cor:interval}), and the exclusion of $-4$ at order $12$, which
gives $13\le\nmin(-4)\le15$ (Theorem~\ref{thm:search}, computer-assisted). The
order-$16$ certificate gives a simple-root example at the remaining order, $16$,
which joins of the $15$-vertex graph with stars do not reach
(Theorem~\ref{thm:orders}). The joining construction incorporates $-10$ into
the prescribed-multiplicity statement (Theorem~\ref{thm:mult}). The edge lists,
verification code and search logs are available at \cite{Art}.

\section{Three lemmas}\label{sec:lemmas}

For a vertex set $S$ write $N[S]$ for the set of vertices in $S$ or adjacent to
a vertex of $S$. For a graph $F$ and $T\subseteq V(F)$ let
\[
  Q_{F,T}(x)=\sum_{\substack{S\subseteq V(F)\\ V(F)\setminus T\subseteq N[S]}} x^{|S|},
  \qquad
  U_{F,T}(x)=\sum_{\substack{S\subseteq V(F)\setminus T\\ S \text{ dominates } F}} x^{|S|}.
\]
Thus $Q_{F,T}$ counts sets that dominate every vertex outside $T$, and
$U_{F,T}$ counts dominating sets disjoint from $T$; $Q_{F,\emptyset}=D(F,x)$.
A \emph{leaf} is a vertex of degree one and a \emph{support} is a vertex
adjacent to a leaf. On vertex sets $\{0,\dots,b-1\}$ with $b\le 10$ we write the
edge $\{u,v\}$ as the two-digit word $uv$.

The first lemma is the case of a cut vertex with pendant leaves of the
articulation-vertex splitting formulas of \cite[Section~4]{KPSTT12}; in this form
it is the template used in \cite{MA26}.

\begin{lemma}[hub construction]\label{lem:hub}
Let $F_1,\dots,F_m$ be disjoint graphs, $T_i\subseteq V(F_i)$, and $\ell\ge 1$.
Let $G$ consist of the $F_i$, a new vertex $h$ (the hub) joined to every
vertex of $T_1\cup\dots\cup T_m$, and $\ell$ new leaves adjacent only to $h$.
Then
\[
  D(G,x)=x^{\ell}\prod_{i=1}^m D(F_i,x)+x(1+x)^{\ell}\prod_{i=1}^m Q_{F_i,T_i}(x).
\]
Consequently, for an integer $k\ge 2$,
\begin{equation}\label{eq:hub}
  D(G,-k)=(-1)^{\ell}k\Bigl(k^{\ell-1}\prod_i D(F_i,-k)-(k-1)^{\ell}\prod_i Q_{F_i,T_i}(-k)\Bigr).
\end{equation}
If every $F_i$ is connected and every $T_i$ is nonempty, $G$ is connected.
\end{lemma}

\begin{proof}
Let $S\subseteq V(G)$ and $S_i=S\cap V(F_i)$. Suppose $h\notin S$. A leaf $p$
has closed neighbourhood $\{p,h\}$, so $S$ dominates $p$ exactly when $p\in S$;
then $h$ is dominated because $\ell\ge1$. A vertex of $F_i$ has no neighbour in
$S$ outside $F_i$, so it is dominated by $S$ exactly when it is dominated by
$S_i$ in $F_i$. Hence these dominating sets contribute
$x^{\ell}\prod_i D(F_i,x)$. Suppose $h\in S$. Then $h$, the leaves and the
vertices of each $T_i$ are dominated, the leaves may or may not lie in $S$,
and a vertex of $F_i$ outside $T_i$ is dominated exactly when it lies in
$N[S_i]$. These sets contribute $x(1+x)^{\ell}\prod_i Q_{F_i,T_i}(x)$.
Substituting $x=-k$ gives $(-k)^{\ell}=(-1)^{\ell}k^{\ell}$ and
$-k(1-k)^{\ell}=-(-1)^{\ell}k(k-1)^{\ell}$, which is \eqref{eq:hub}. The last
sentence is clear.
\end{proof}

\begin{lemma}[joining supports; {\cite[Theorems~3.10, 3.11]{KPSTT12}}]\label{lem:join}
Let $G_1,\dots,G_r$ be disjoint graphs and $s_i$ a support of $G_i$. If $G$ is
obtained from $G_1\cup\dots\cup G_r$ by adding any set of edges between the
vertices $s_1,\dots,s_r$, then $D(G,x)=\prod_{i=1}^r D(G_i,x)$.
\end{lemma}

\begin{proof}
Adding edges cannot destroy domination, so if each $S\cap V(G_i)$ dominates
$G_i$ then $S$ dominates $G$. Conversely let $S$ dominate $G$, put
$S_i=S\cap V(G_i)$, and let $p_i$ be a leaf of $G_i$ adjacent to $s_i$. The new
edges do not meet $p_i$, so $p_i$ has closed neighbourhood $\{p_i,s_i\}$ in $G$
and $S$ contains $p_i$ or $s_i$; either way $s_i\in N[S_i]$. Every other vertex
of $G_i$ has the same closed neighbourhood in $G$ as in $G_i$. Hence $S_i$
dominates $G_i$. The dominating sets of $G$ are therefore exactly the unions of
dominating sets of the $G_i$, and $x^{|S|}=\prod_i x^{|S_i|}$.
\end{proof}

A support is domination-covered because its leaf has closed neighbourhood
contained in that of the support. By Kotek et al.\ \cite[Theorem~3.11]{KPSTT12},
an edge $uv$ is irrelevant in $G$ if its endpoints are domination-covered in
$G-uv$. In particular, an edge can be added between nonadjacent supports without
changing the polynomial. The same lemma is used in \cite{MA26} and \cite{MZ26}.

The next lemma writes the vertex-inclusion and exclusion split of Kotek et
al.\ \cite[Theorem~2.1]{KPSTT12} in the form used by the exhaustive search of
Section~\ref{sec:four}.

\begin{lemma}[one more vertex]\label{lem:ext}
Let $H$ be a graph, $\emptyset\ne T\subseteq V(H)$, and let $H+v_T$ be $H$
together with a new vertex $v$ adjacent exactly to $T$. Then
\[
  D(H+v_T,x)=D(H,x)-U_{H,T}(x)+x\,Q_{H,T}(x).
\]
\end{lemma}

\begin{proof}
A set $S\subseteq V(H)$ dominates $H+v_T$ exactly when it dominates $H$ and
meets $T$; these sets contribute $D(H,x)-U_{H,T}(x)$. A set $S\cup\{v\}$ with
$S\subseteq V(H)$ dominates $H+v_T$ exactly when $S$ dominates every vertex of
$H$ outside $T$; these contribute $x\,Q_{H,T}(x)$.
\end{proof}

\section{A domination root at \texorpdfstring{$-10$}{-10}}\label{sec:ten}

Let $F_1,\dots,F_{\TenTypes}$ be the connected graphs listed in
Appendix~\ref{app:blocks}, with attachment sets $T_i$ and multiplicities $c_i$
as in Table~\ref{tab:blocks}. The graph $G_{10}$ is the graph of
Lemma~\ref{lem:hub} with $\ell=\TenLeaves$ leaves and $c_i$ disjoint copies of
each $F_i$, the hub being joined to the copy of $T_i$ in every copy of $F_i$.
So $G_{10}$ is a vertex with $\TenLeaves$ pendant leaves and $\TenBlocks$
attached blocks of at most $10$ vertices each. It has
$1+\TenLeaves+\sum_i c_i|V(F_i)|=\TenOrder$ vertices and
$\TenLeaves+\sum_i c_i(|E(F_i)|+|T_i|)=\TenEdges$ edges.

\begin{table}[ht]
\centering\small
\caption{The blocks of $G_{10}$: number of copies $c_i$, order $b_i$, and the
values at $x=-10$ factored into primes. Edges and attachment sets $T_i$ are in
Table~\ref{tab:edges}.}
\label{tab:blocks}
\begin{tabular}{rrrll}
\toprule
$i$ & $c_i$ & $b_i$ & $D(F_i,-10)$ & $Q_{F_i,T_i}(-10)$\\
\midrule
\BlockRows
\bottomrule
\end{tabular}
\end{table}

\begin{theorem}\label{thm:ten}
The graph $G_{10}$ is connected and $-10$ is a simple root of $D(G_{10},x)$.
\end{theorem}

\begin{proof}
$G_{10}$ is connected by Lemma~\ref{lem:hub}, since every $F_i$ is connected
and every $T_i$ is nonempty. By \eqref{eq:hub} with $k=10$ and $\ell=\TenLeaves$,
$D(G_{10},-10)=0$ if and only if
\begin{equation}\label{eq:ten}
  10^{21}\prod_{i=1}^{\TenTypes} D(F_i,-10)^{c_i}
  \;=\;9^{22}\prod_{i=1}^{\TenTypes} Q_{F_i,T_i}(-10)^{c_i}.
\end{equation}
The entries of Table~\ref{tab:blocks} are computed from the definitions by
listing the at most $2^{10}$ subsets of $V(F_i)$. In each row the two values
have the same sign, so both sides of \eqref{eq:ten} have the same sign. Reading
prime exponents off the table, with multiplicities $c_i$: the exponent of $3$ is
$\TenThree$ on both sides (every left entry except row~3 contributes one factor
$3$, and $\sum_{i\ne 3}c_i=44$); the exponent of $2$ is $\TenTwo$ and the
exponent of $5$ is $\TenFive$ on both sides; and each other prime appearing
in Table~\ref{tab:blocks} has the same total exponent on the two sides. Hence
\eqref{eq:ten} holds and $D(G_{10},-10)=0$. By Table~\ref{tab:residues},
$D'(G_{10},-10)\equiv 912\pmod{1009}$, so the root is simple.
\end{proof}

\begin{remark}
The full polynomial $D(G_{10},x)$, computed in two independent ways
(Section~\ref{sec:certificates}), is
$x^{\TenGamma}(x+2)^{16}(x+10)R(x)$ with $R(r)\ne 0$ for every integer
$r\in[-11,-1]$. The search at $-10$ uses the block-ratio catalog and
prime-valuation approach that MathArena used at $-6$ and $-8$ \cite{MA26}. Its
catalog had the ratios $D(F,-10)/Q_{F,T}(-10)$ for every connected graph $F$
with at most $8$ vertices and for random connected graphs of larger order,
with every attachment set $T$. It retained nonzero ratios whose reductions use
only primes at most $97$, then an integer program balanced their valuations,
keeping the order small, and the signs in \eqref{eq:ten} were checked. No
minimality is claimed.
\end{remark}

\begin{corollary}\label{cor:interval}
Let $r$ be an integer with $-11\le r\le -1$. Some graph has $D(G,r)=0$ if and
only if $r\in\{-2,-4,-6,-8,-10\}$, and in that case $G$ can be taken connected.
\end{corollary}

\begin{proof}
For odd $r$, $D(G,r)\equiv D(G,1)\equiv 1\pmod 2$ by \cite[Theorem~1.1]{BCS09}.
Connected examples for the even values are $K_2$ for $-2$, the graph of
\cite{AG26} (or the smaller ones of \cite{MA26,Om26,Res26}) for $-4$, the
graphs of \cite{MA26} for $-6$ and $-8$, and $G_{10}$ for $-10$.
\end{proof}

\section{Small graphs with domination root \texorpdfstring{$-4$}{-4}}\label{sec:four}

\begin{theorem}\label{thm:search}
No graph with at most $12$ vertices has domination root $-4$. Hence
$13\le \nmin(-4)\le 15$.
\end{theorem}

\begin{proof}[Proof (computer-assisted)]
Because $D$ is multiplicative over components it suffices to consider connected
graphs. The upper bound is the $15$-vertex graph of \cite{MA26}.

\emph{Orders $1$ to $11$.} The program \texttt{geng} of nauty~2.8.9
\cite{MP14} lists one graph from each isomorphism class of connected graphs of
a given order. For every listed graph the program \texttt{search/domproc.c}
counts the dominating sets of each size by running through all $2^n$ vertex
subsets, and evaluates $D(G,r)$ for $r=-2,\dots,-10$ in exact $128$-bit integer
arithmetic. The numbers of graphs listed for $n=1,\dots,11$ are $1$, $1$, $2$,
$6$, $21$, $112$, $853$, $11117$, $261080$, $11716571$, $1006700565$, in
agreement with OEIS A001349 \cite{OEIS}; no value $D(G,-4)$ is zero.

\emph{Order $12$.} Every connected graph $G$ on $12$ vertices has a vertex $v$
with $G-v$ connected (a leaf of a spanning tree), so $G\cong H+v_T$ for a
connected graph $H$ on $11$ vertices and some nonempty $T\subseteq V(H)$, in the
notation of Lemma~\ref{lem:ext}; relabelling, we may take $H$ to be the
representative of its isomorphism class listed by \texttt{geng}. For every one of the $1006700565$ graphs $H$
listed by \texttt{geng}, the program \texttt{search/extproc.c} computes
$D(H,-4)$, $U_{H,T}(-4)$ and $Q_{H,T}(-4)$ for all $2^{11}$ sets $T$ at once.
It adds the weight $(-4)^{|S|}$ of every $S\subseteq V(H)$ to entry
$V(H)\setminus N[S]$ of a first array and, when $S$ dominates $H$, to entry $S$
of a second array. After a subset-sum transform of both arrays (each entry
replaced by the sum of the entries indexed by its subsets; MathArena uses a
subset transform for its attachment counts \cite{MA26}), entry $T$ of the
first array is $Q_{H,T}(-4)$ and entry $V(H)\setminus T$ of the second is
$U_{H,T}(-4)$. Every array entry has absolute value at most
$\sum_S 4^{|S|}=5^{11}$ and every value of Lemma~\ref{lem:ext} at most
$6\cdot 5^{11}<2^{63}$, so $64$-bit arithmetic is exact. All
$2047\cdot 1006700565=2060716056555$ values of $D(H+v_T,-4)$ given by
Lemma~\ref{lem:ext} are nonzero.
\end{proof}

\begin{remark}\label{rem:others}
The same computation shows that no graph with at most $11$ vertices has a
domination root at $-6$, $-8$ or $-10$, so $12\le\nmin(-6)\le 31$,
$12\le\nmin(-8)\le 83$ (upper bounds from \cite{MA26}) and
$12\le\nmin(-10)\le 373$.
\end{remark}

MathArena reports witnesses of orders $15$ and $16$ \cite{MA26}; with its join
remark and stars these already give connected graphs with root $-4$ of every
order $n\ge15$. The next theorem writes this out, with published data at order
$16$ and a simple root at every order. Let $F_{16}$ be the connected graph on
$\{0,\dots,9\}$ with the $16$ edges
\[
  \SixteenEdges,
\]
and let $G_{16}$ be the graph of Lemma~\ref{lem:hub} with this single block, the
nonempty attachment set $T=\SixteenT$ and $\ell=5$ leaves. By
Lemma~\ref{lem:hub} it is connected; it has $16$ vertices and $16+4+5=25$ edges.

\begin{theorem}[cf.~\cite{MA26}]\label{thm:orders}
$D(G_{16},-4)=0$, and the root $-4$ is simple. For every $n\ge 15$ there is a
connected graph of order $n$ in which $-4$ is a simple domination root.
\end{theorem}

\begin{proof}
Listing the $2^{10}$ subsets of $V(F_{16})$ gives $D(F_{16},-4)=\SixteenP$ and
$Q_{F_{16},T}(-4)=\SixteenQ$. Since
$4^4\cdot\SixteenP=3^5\cdot\SixteenQ=\SixteenProd$, \eqref{eq:hub} with $k=4$ and
$\ell=5$ gives $D(G_{16},-4)=0$. Direct enumeration of the $2^{16}$ subsets
gives $D(G_{16},-4)=0$ and $D'(G_{16},-4)=\SixteenDeriv$, so the root is
simple.

For $n=15$ take the graph $G_4$ of \cite{MA26}: on $\{0,\dots,14\}$ with edges
\FourEdges; it has
$D(G_4,-4)=0$ and $D'(G_4,-4)=\FourDeriv$. For $n=16$ take $G_{16}$. For
$n\ge17$ let $t=n-16\ge1$ and join the support $2$ of $G_4$ (adjacent to the
leaves $10,\dots,14$) by one edge to the centre of the star $K_{1,t}$, which is
a support. The result is connected of order $n$, and by
Lemma~\ref{lem:join} its domination polynomial is $D(G_4,x)D(K_{1,t},x)$, where
$D(K_{1,t},x)=x(1+x)^t+x^t$. Now
$D(K_{1,t},-4)=(-1)^t\,4\,(4^{t-1}-3^t)\ne0$, because $4^{t-1}-3^t$ is $-2$
for $t=1$ and odd for $t\ge2$. So $-4$ is a simple root of the product.
\end{proof}

\begin{remark}
$G_{16}$ was found by a computer search among one-vertex extensions of
$15$-vertex graphs with root $-4$ that have the shape of Lemma~\ref{lem:hub}
with a single block. $G_4$ itself has that shape: the hub $2$, five leaves, and a $9$-vertex
block attached at $\{3,5\}$.
\end{remark}

\section{Prescribed multiplicities}\label{sec:mult}

Let $G_4$, $G_6$, $G_8$ be the connected graphs of \cite{MA26} with roots
$-4$, $-6$, $-8$ and $15$, $31$, $83$ vertices, and let $G_{10}$ be as in
Section~\ref{sec:ten}. Each has a support: the centre of $G_4$, $G_6$, $G_8$
(adjacent to $5$, $10$ and $16$ leaves) and the hub of $G_{10}$.

\begin{table}[ht]
\centering\small
\caption{Residues modulo $1009$ of $D(G_k,-j)$ and of $D'(G_k,-k)$. The
diagonal values $D(G_k,-k)$ are exactly $0$; every other residue is nonzero,
hence so is the integer.}
\label{tab:residues}
\begin{tabular}{lrrrrrr}
\toprule
 & order & $D(G_k,-4)$ & $D(G_k,-6)$ & $D(G_k,-8)$ & $D(G_k,-10)$ & $D'(G_k,-k)$\\
\midrule
\ResidueRows
\bottomrule
\end{tabular}
\end{table}

\begin{theorem}\label{thm:mult}
For all nonnegative integers $m_4$, $m_6$, $m_8$ and $m_{10}$ there is a
connected graph $G$ such that the multiplicity of $x+k$ as a factor of
$D(G,x)$ is $m_k$ for each $k\in\{4,6,8,10\}$.
\end{theorem}

\begin{proof}
If all $m_k$ are $0$, take $K_1$, with $D(K_1,x)=x$. Otherwise take $m_k$
disjoint copies of $G_k$ for each $k$ and join their designated supports along
a path. The result is connected, and by Lemma~\ref{lem:join} its domination
polynomial is $\prod_k D(G_k,x)^{m_k}$. By Table~\ref{tab:residues},
$-k$ is a simple root of $D(G_k,x)$ and is not a root of $D(G_j,x)$ for
$j\ne k$; so its multiplicity in the product is $m_k$.
\end{proof}

MathArena's join and the off-diagonal checks in Table~\ref{tab:residues} give
prescribed multiplicities for its three roots \cite{MA26}. Ma and Zhang give
exact multiplicities for $-4$ alone \cite{MZ26}. The $-10$ witness extends the
construction to four roots.

\section{Computational certificates}\label{sec:certificates}

The program \texttt{verify.py} uses only the Python standard library. It reads
the graphs from \texttt{data/graphs.json} and \texttt{data/G10-edges.txt},
recomputes everything from the definitions, regenerates the numbers in
Tables~\ref{tab:blocks}, \ref{tab:residues}, \ref{tab:edges} and in
Theorem~\ref{thm:orders} and checks them against the file included by this
paper, and stops at the first discrepancy. Its default mode runs in about one
second and checks:
\begin{itemize}
\item Theorem~\ref{thm:ten}: the edge list is the graph described by the block
  data; $G_{10}$ is simple and connected with $\TenOrder$ vertices and
  $\TenEdges$ edges; $D(G_{10},x)$ is computed from the edge list alone by a
  dynamic program over a vertex order (each processed vertex that still has an
  unprocessed neighbour carries one of three states: in $S$, dominated, not yet
  dominated), and again from Lemma~\ref{lem:hub} with each block counted by
  subset enumeration; the two polynomials are equal; the multiplicities of its
  integer roots in $[-11,-1]$; and identity \eqref{eq:ten}. The dynamic program
  is first checked against subset enumeration on $300$ random graphs.
\item Theorem~\ref{thm:search} at small scale: a separate Python enumeration of
  connected graphs up to isomorphism through order $7$ (counts as in
  A001349), no root at $-4,-6,-8,-10$ there, and no root $-4$ at order $8$ via
  Lemma~\ref{lem:ext} and the subset transform. For all extensions of graphs
  of order at most $5$, Lemma~\ref{lem:ext} is compared with subset
  enumeration, and every value of the subset transform at $x=-2$ and $x=-4$ with
  the directly computed polynomial; this includes extensions with a root at
  $-2$, so the transform is tested on hits as well as on nonzero values.
\item Theorem~\ref{thm:orders}: $G_{16}$ and $G_4$ by enumerating all subsets;
  the joins with $K_{1,t}$ for $t=1,\dots,4$ (orders $17$ to $20$). The full
  polynomial $D(G_{16},x)$ is the macro \texttt{\textbackslash SixteenPoly} of
  the generated file \texttt{tables.tex}.
\item Theorem~\ref{thm:mult}: Table~\ref{tab:residues} from the exact
  polynomials, and, for the multiplicities $m_4=2$ and $m_6=m_8=m_{10}=1$, the
  joined graph on $517$ vertices: its polynomial, computed by the dynamic
  program, equals the product and has the stated multiplicities.
\end{itemize}
The command
\begin{center}
\texttt{python3 -B verify.py --full NAUTY\_DIR}
\end{center}
also recompiles \texttt{search/domproc.c} and \texttt{search/extproc.c} against
a built copy of nauty~2.8.9 and reruns the search of Theorem~\ref{thm:search},
comparing graph counts with A001349 and, for orders up to $7$, the number of
graphs with root $-2$ with the Python enumeration. A fresh run compiles both
processors and executes every partition. A resumed run reuses a partition only
when its program fingerprint, arguments and complete result record match.
Every record, fresh or reused, must consist of one result line with the
expected fields, order and root and nonnegative counts, and for order $12$ the
number of tests in each partition must be $2047$ times its number of base
graphs. Before the search, \texttt{search/extproc.c} compiled on its own must
report the root $-4$ of $G_4$ from $G_4$ minus a leaf, and agree with
Lemma~\ref{lem:ext} on that graph. On an AMD Ryzen~9 4900H (8 cores, 16
threads) with $15$ parallel jobs, order $11$ took $6$ to $9$ minutes and order
$12$ about an hour ($3518$ seconds in the first run). The direct search at order
$10$ took about $18$ seconds in the first run. The search was run twice, with
identical results; the logs, the build commands and the SHA-256 hashes of the
sources and binaries are in \texttt{logs/}.

\section{Open questions}

\begin{enumerate}
\item Is every negative even integer a domination root of some graph? The
  first case not settled here is $-12$.
\item Is $\nmin(-4)$ equal to $13$, $14$ or $15$? Extending the search of
  Theorem~\ref{thm:search} to order $13$ means testing all nonempty
  neighbourhoods of a new vertex for each of the $164059830476$ connected
  graphs on $12$ vertices \cite{OEIS}: about $160$ times as many base graphs
  as the order-$12$ run, each with twice as many sets $T$.
\item What are $\nmin(-6)$, $\nmin(-8)$ and $\nmin(-10)$? By
  Remark~\ref{rem:others} they lie in $[12,31]$, $[12,83]$ and $[12,373]$.
\end{enumerate}

\section*{Use of AI and verification}

The results were found with AI systems directed by the author. The language
model GPT-6.2 Astra, working with computer search, found $G_{10}$ and wrote the
search that produced it. Anthropic's Claude Opus~5.5 rechecked the claims with
new code, searched the literature, wrote \texttt{search/} and
\texttt{verify.py}, ran the searches of Theorem~\ref{thm:search}, found
$G_{16}$ and drafted this text. The checks are those of
Section~\ref{sec:certificates}: two independent computations of
$D(G_{10},x)$, subset enumeration of every block and of $G_{16}$ and $G_4$, a
separate Python enumeration of small graphs, and two runs of the exhaustive
search with identical results. No AI system is an author; the author is
responsible for the content.

The computer-assisted parts are: the entries of Tables~\ref{tab:blocks} and
\ref{tab:residues}, the values $D(F_{16},-4)$, $Q_{F_{16},T}(-4)$,
$D'(G_{16},-4)$ and $D'(G_4,-4)$, the remark on the factorisation of
$D(G_{10},x)$ (all reproduced by \texttt{python3 -B verify.py}), and the whole
of Theorem~\ref{thm:search} and Remark~\ref{rem:others} (reproduced by
\texttt{python3 -B verify.py --full NAUTY\_DIR}). Lemmas~\ref{lem:hub}--\ref{lem:ext},
Corollary~\ref{cor:interval} and the deductions in Theorems~\ref{thm:ten},
\ref{thm:orders} and~\ref{thm:mult} are proved by hand above.

\begin{thebibliography}{99}

\bibitem{AAOP10}
S.~Akbari, S.~Alikhani, M.~R.~Oboudi and Y.-H.~Peng,
\emph{On the zeros of domination polynomial of a graph},
Contemporary Mathematics \textbf{531} (2010), 109--115.
\url{https://doi.org/10.1090/conm/531/10460}

\bibitem{AG26}
S.~Alikhani and M.~Griswold,
\emph{On the integer domination root conjecture},
arXiv:2608.00109v1 [math.CO], 31 July 2026.
\url{https://arxiv.org/abs/2608.00109}

\bibitem{Art}
A.~Poonacha, \emph{Data, code and search logs for ``A graph with domination
root $-10$, and small graphs with domination root $-4$''}, release
\texttt{domination-roots-v1} of the repository \texttt{math-findings}.
\url{https://github.com/Aryan-Poonacha/math-findings/tree/domination-roots-v1/domination-roots}

\bibitem{BCS09}
A.~E.~Brouwer, P.~Csorba and A.~Schrijver,
\emph{The number of dominating sets of a finite graph is odd},
preprint, 2 June 2009.
\url{https://www.win.tue.nl/~aeb/preprints/domin4a.pdf} (accessed 27 November 2026)

\bibitem{KPSTT12}
T.~Kotek, J.~Preen, F.~Simon, P.~Tittmann and M.~Trinks,
\emph{Recurrence relations and splitting formulas for the domination polynomial},
Electron. J. Combin. \textbf{19}(3) (2012), \#P47.
\url{https://doi.org/10.37236/2475}

\bibitem{MZ26}
H.~Ma and W.~Zhang,
\emph{Alikhani--Griswold, Section~4(3): connected graphs with domination root $-4$},
Zenodo, 3 November 2026.
\url{https://doi.org/10.5281/zenodo.23919066}

\bibitem{MA26}
MathArena,
\emph{The counterexample hiding in a green cell} (model submission, GPT-6.1 Sol),
blog post with witness data.
\url{https://matharena.ai/static/blog_submissions/llm_blog_posts/gpt-61-sol/blog.html};
archived 10 October 2026 at
\url{https://web.archive.org/web/20261010115259/https://matharena.ai/static/blog_submissions/llm_blog_posts/gpt-61-sol/blog.html}
(accessed 27 November 2026)

\bibitem{MP14}
B.~D.~McKay and A.~Piperno,
\emph{Practical graph isomorphism, II},
J. Symbolic Comput. \textbf{60} (2014), 94--112.
\url{https://doi.org/10.1016/j.jsc.2013.09.003};
software at \url{https://pallini.di.uniroma1.it/} (accessed 27 November 2026)

\bibitem{OEIS}
OEIS Foundation, \emph{Sequence A001349: number of connected simple graphs with
$n$ nodes}, The On-Line Encyclopedia of Integer Sequences.
\url{https://oeis.org/A001349} (accessed 27 November 2026)

\bibitem{Om26}
The Omega Institute, \emph{trureturing}, Lean declaration
\texttt{AlikhaniGriswoldMinOrderRefutation.result} (a $21$-vertex graph with
domination root $-4$), 1 November 2026.
\url{https://github.com/the-omega-institute/trureturing/issues/25239}
(accessed 27 November 2026)

\bibitem{Om26b}
The Omega Institute, \emph{Complete the connected domination-root minus-four
infinite family}, pull request \#25225 of \emph{trureturing}, merged 1 November
2026. \url{https://github.com/the-omega-institute/trureturing/pull/25225}
(accessed 27 November 2026)

\bibitem{Res26}
V.~Reshetnikov (ed.),
\emph{A bicyclic 30-vertex graph with domination root $-4$}, Part~I
(18 September 2026) and Part~II (30 September 2026), report in the ProveIt
repository.
\url{https://github.com/VladimirReshetnikov/ProveIt/tree/main/SetTheory/Cardinals/docs/reports/graph-theory/domination-root-minus-four-order-30}
(accessed 27 November 2026)

\bibitem{Wei}
E.~W.~Weisstein, \emph{Domination root}, MathWorld---A Wolfram Web Resource.
\url{https://mathworld.wolfram.com/DominationRoot.html} (accessed 27 November 2026)

\end{thebibliography}

\appendix
\section{The blocks of \texorpdfstring{$G_{10}$}{G10}}\label{app:blocks}

Block $F_i$ has vertex set $\{0,\dots,b_i-1\}$, and its edges are written as
two-digit words (Section~\ref{sec:lemmas}). Multiplicities are in Table~\ref{tab:blocks}. Blocks $10$ and $11$ are the same graph with different
attachment sets. The same data are in \texttt{data/graphs.json}, and the
expanded edge list of $G_{10}$ (hub $0$, leaves $1$--$22$, then the block copies
in order) is \texttt{data/G10-edges.txt}.

\begin{table}[ht]
\centering\small
\caption{Attachment sets and edges of the blocks $F_1,\dots,F_{\TenTypes}$.}
\label{tab:edges}
\begin{tabular}{rll}
\toprule
$i$ & $T_i$ & edges of $F_i$\\
\midrule
\BlockEdgeRows
\bottomrule
\end{tabular}
\end{table}

\end{document}
